\section{\label{sec:background}Background}

\subsection{Electromagnetic simulation of photonic devices}

We consider two-dimensional TE-polarized electromagnetic fields in silicon-on-insulator (SOI) waveguides.
The out-of-plane electric field component $E_z(x,y)$ satisfies the scalar Helmholtz equation in the frequency domain:
\begin{equation}
    \nabla^2 E_z(x,y) + k_0^2 \,\varepsilon(x,y)\, E_z(x,y) = 0,
    \label{eq:helmholtz}
\end{equation}
where $k_0 = 2\pi/\lambda$ is the free-space wavenumber and $\varepsilon(x,y)$ is the spatially varying relative
permittivity. The reduction from three dimensions to two is valid when the vertical slab mode has been pre-solved: for
a 220\,nm SOI slab the fundamental TE mode yields an effective core index $n_\text{eff}\approx 2.4$
($\varepsilon_\text{core}\approx 5.8$) and the silica cladding has $n_\text{clad}=1.444$
($\varepsilon_\text{clad}\approx 2.09$)~\cite{Reed2004SiliconPhotonics}. We operate in the telecommunications C-band,
$\lambda\in[1.50,1.60]\,\mu$m.

Finite-difference time-domain (FDTD) methods solve the time-domain Maxwell's equations on a Yee grid, exciting the
structure with an eigenmode source and extracting frequency-domain fields via discrete Fourier
transforms~\cite{TafloveHagness2005FDTD}. We use the open-source Meep solver~\cite{Oskooi2010Meep}. Each simulation runs until field energy decays below a convergence
threshold, typically requiring minutes for the domain sizes considered here.

The key outputs of an FDTD simulation are the complex field $E_z(x,y)$, the permittivity map $\varepsilon(x,y)$, and
the scattering parameters (S-parameters) $S_{ji}\in\mathbb{C}$ that characterize port-to-port transfer. The
S-parameter from input port $i$ to output port $j$ is defined as the ratio of the outgoing mode amplitude at port~$j$
to the incident mode amplitude at port~$i$:
\begin{equation}
    S_{ji} = \frac{a_j}{a_i},
    \label{eq:sparam_def}
\end{equation}
where the mode amplitudes $a_p$ are obtained by projecting the total field onto the fundamental waveguide eigenmode at
each port cross-section. The squared magnitude $|S_{ji}|^2$ gives the fractional power transfer; for example, a
directional coupler is characterized by its splitting ratio $|S_{31}|^2/|S_{41}|^2$ and a waveguide bend by its
insertion loss $-10\log_{10}|S_{21}|^2$.  S-parameters are the primary design metric for photonic components and the
quantities that must ultimately be predicted accurately by any surrogate model.

\subsection{Conditional flow matching}

Flow matching~\cite{Lipman2022FlowMatching,Liu2023RectifiedFlow} is a generative modeling framework that learns a
time-dependent velocity field $u_\theta(x_t,t)$ to transport samples from a simple prior $p_0$ (Gaussian noise) to the
data distribution $p_1$ along a deterministic path. Given source--target pairs $x_0\sim\mathcal{N}(0,I)$ and
$x_1\sim p_\text{data}$, the conditional interpolation path is
\begin{equation}
    x_t = (1-t)\,x_0 + t\,x_1, \quad t\in[0,1],
    \label{eq:interpolation}
\end{equation}
with corresponding target velocity $u_t = x_1 - x_0$. A neural network $u_\theta$ is trained to match this velocity
by minimizing
\begin{equation}
    \mathcal{L}_\text{FM} = \mathbb{E}_{t,\,x_0,\,x_1}\!\left[\left\|u_\theta(x_t,t) - u_t\right\|^2\right].
    \label{eq:fm_loss}
\end{equation}
At inference, new samples are generated by integrating the learned velocity field from $t{=}0$ to $t{=}1$ using an ODE
solver, starting from Gaussian noise. Compared with score-based diffusion models~\cite{Song2021ScoreSDE}, flow matching yields straighter transport paths
that require fewer integration steps and avoids the need to tune a noise
schedule~\cite{Lipman2022FlowMatching,Liu2023RectifiedFlow}.

In our setting, the data distribution $p_1$ consists of physically valid electromagnetic field solutions conditioned on
device geometry and wavelength.  The model learns to denoise a random field into one that satisfies the Helmholtz
equation for a given permittivity map --- a task we augment with explicit physics constraints, described in
Sec.~\ref{sec:method}.
